% ============================================================
%  Front matter — capa, falsa-capa, dedicatória, epígrafe,
%  resumo, agradecimentos, abreviaturas
% ============================================================

% ---------- Página de rosto ----------
\begin{titlingpage}
\thispagestyle{empty}
\begin{center}
\vspace*{3cm}
{\Huge\scshape História do Inferno\par}
\vspace{0.8cm}
{\Large\slshape Edição Expandida e Complementada\par}
\vspace{2.5cm}
{\large Uma Análise Interdisciplinar dos Conceitos de\\
       Punição Eterna da Pré-História à Era Digital\par}
\vspace{2.0cm}
{\large Baseada em 195 Fontes Acadêmicas Especializadas\par}
\vfill
{\large Ney Lemke\par}
\vspace{1.5cm}
{\small São Paulo --- Botucatu\par}
{\small \the\year\par}
\end{center}
\end{titlingpage}

% ---------- Verso da página de rosto (nota editorial) ----------
\thispagestyle{empty}
\vspace*{2cm}
\begin{quote}
\small
\noindent\textbf{Nota editorial sobre esta edição.}

Esta obra foi convertida automaticamente do formato \texttt{.docx}
para \LaTeX{} (\textit{memoir}, classe de livro) preservando a
hierarquia original de 5 partes e 20 capítulos, além de 83 imagens
extraídas para o diretório \texttt{figuras/}.

Durante a conversão foram identificadas duas inconsistências no
original: (i) o sumário referencia o \emph{Capítulo 4 (Período de
Transição Antiga)}, ausente do corpo do texto; (ii) o \emph{Capítulo
17} aparece no corpo sob dois títulos divergentes (``O Inferno no
divã'' e ``O Inferno Contemporâneo''). A numeração impressa segue o
corpo do manuscrito.
\end{quote}
\vfill
\clearpage

% ---------- Dedicatória ----------
\thispagestyle{empty}
\vspace*{4cm}
\begin{center}
\itshape
``A questão do inferno não é apenas sobre o que nos espera depois da
morte.\\[0.3cm]
É sobre o que somos capazes de fazer \emph{aqui}.''
\end{center}
\clearpage

% ---------- Epígrafe geral ----------
\thispagestyle{empty}
\vspace*{3cm}
\begin{flushright}
\begin{minipage}{0.7\textwidth}
\itshape
Há mais coisas entre o céu e a terra, Horacio,\\
do que sonha a nossa vã filosofia.
\end{minipage}

\hfill --- Shakespeare, \textit{Hamlet}, Ato I, Cena V
\end{flushright}
\vfill
\clearpage

% ---------- Resumo ----------
\chapter*{Resumo}
\addcontentsline{toc}{chapter}{Resumo}

Esta obra apresenta uma análise abrangente e interdisciplinar da
evolução histórica dos conceitos de inferno e punição eterna, desde
as primeiras práticas funerárias da pré-história até as
representações distópicas da era digital. Baseada em 195 fontes
acadêmicas especializadas, a investigação percorre cinco grandes
eixos: (i) as origens e a Antiguidade Clássica, com atenção ao
submundo sumério, egípcio e greco-romano; (ii) as tradições
regionais e culturais, incluindo as mitologias nórdica, celta,
eslava, bizantina e ameríndia; (iii) os desenvolvimentos judaico-cristãos,
do \emph{Sheol} à patrística; (iv) as sínteses medievais e modernas,
do Islã medieval a Dante, da Renascença ao Iluminismo; e (v) a era
contemporânea, com suas manifestações literárias, cinematográficas,
digitais e de gênero.

A tese central sustenta que a ideia de inferno --- longe de ser
monolítica --- é uma construção cultural cumulativa, reelaborada a
cada época conforme as ansiedades e os dil éticos de cada civilização.
A análise dialoga com a história das religiões, a teologia, a
antropologia, a literatura, o cinema e os estudos culturais.

% ---------- Agradecimentos ----------
\chapter*{Agradecimentos}
\addcontentsline{toc}{chapter}{Agradecimentos}

Esta obra resulta de anos de investigação em fontes dispersas por
bibliotecas, manuscritos digitalizados, arquivos e bases de dados
acadêmicas. Agradeço a todos os pesquisadores, tradutores e
bibliotecários que, anonimamente, mantêm viva a memória documental
da humanidade. Em especial, aos colegas da UNESP que, em conversas
informais, ajudaram a amadurecer as hipóteses aqui apresentadas.

% ---------- Lista de abreviaturas (opcional, expandível) ----------
\chapter*{Nota sobre abreviaturas e transliterações}
\addcontentsline{toc}{chapter}{Nota sobre abreviaturas e transliterações}

As transliterações de termos em sânscrito, acádio, hebraico, árabe e
grego seguem convenções padronizadas da literatura acadêmica de
referência. Termos em latim aparecem em \textit{itálico} quando
conservados na forma original; termos gregos seguem o sistema
betacode latino. Para facilitar a consulta, o índice remissivo ao
final do volume remete o leitor tanto à grafia original quanto à
tradução aportuguesada.
\clearpage